% output=pdftex interface=nl 

\environment man-init \hoofdtaal[nl]

\starttekst

\titelblad
  {PPCH\TeX}
  {een serie macro's voor het zetten\\ 
   van chemische structuur formules}
  {J. Hagen \& A.F. Otten\\
   Pragma ADE, Hasselt NL}

\titel{Inhoud}

\stelpaginanummerin[status=stop] \plaatsinhoud \pagina

\titel{Inleiding}

\PPCHTEX\ is een samenhangende serie macro's waarmee chemische formules
kunnen worden gezet. De macro's vallen terug op \PICTEX, een in een door
Michael Wichura in public domain gebracht macropakket, waarmee grafieken en
andere lijnafbeeldingen kunnen worden getekend. Hoewel \PPCHTEX\ in eerste
instantie is ontwikkeld voor \PICTEX, kan inmiddels ook gebruik worden
gemaakt van \PSTRICKS\ van Timothy Van Zandt, zij het met enige beperkingen.
Overigens gebiedt de eerlijkheid te zeggen dat \PICTEX\ de mooiste resultaten
biedt. 

De macro's zijn te gebruiken binnen verschillende \TEX||omgevingen. Daarbij
wordt gebruik gemaakt van enkele \CONTEXT||modules. De macro's zijn zodanig
opgezet dat uitbreiden relatief eenvoudig is. De interactie sluit aan op de
binnen \CONTEXT\ gebruikte interactie. 

Hoewel \PPCHTEX\ in eerste instantie is bedoeld om chemische
structuurformules, zoals zesringen, te zetten, kunnen ook reactiemechanismen
worden weergegeven. De chemische structuren kunnen in verschillende formaten
worden gezet, waarbij vergelijkbare formules optisch op elkaar aansluiten.
Veel voorkomende structuren kunnen worden voorgedefinieerd en opgeroepen. 

Bij de ontwikkeling van de macro's is snelheid ondergeschikt gemaakt aan
flexibiliteit, eenvoud en kwaliteit. Er is geen gebruik gemaakt van het
binnen \PICTEX\ beschikbare mechanisme om delen van figuren op te slaan in
een file. Gebleken is dat daarmee nauwelijks tijdwinst wordt geboekt. Het is
niet anders. 

De eerste versie van \PPCHTEX\ kwam beschikbaar in 1995. Deze handleiding
beschrijft de tweede versie, die dankzij de vele suggesties van Tobias
Burnus, Dirk Kuypers en Ton Otten een vrijwel dubbele functionaliteit kent. 

{\em De source van \PPCHTEX\ is op dit moment geen schoolvoorbeeld van goed
gedocumenteerde code. Te zijner tijd zullen de macro's nog eens onder de loep
worden genomen en gedocumenteerd.} 

\pagina \stelpaginanummerin[status=start]

\vervolgblad
  {Uitleg}
  []
  [SIX,B,C,R1356,RZ1356]
  [NO_2,NO_2,NO_2,CH_3]

\hoofdstuk{Structuren}

Het aantal commando's dat wordt gebruikt om chemische structuurformules te
zetten is, afgezien van wat macro's om toeters en bellen toe te voegen,
beperkt tot vier.\voetnoot{Het begrip structuur heeft in deze handleiding
alleen betrekking op de chemische structuur en niet op de structuur van de
tekst waarin deze formule wordt gezet.} In het volgende voorbeeld zijn al
deze commando's gebruikt: 

\startbuffer
\stelchemiein[assenstelsel=aan,kader=aan]

\startchemie
  \chemie[SIX,B,R,RZ][1,2,3,4,5,6]
\stopchemie
\stopbuffer

\startfiguurtekst
  [links][]
  {geen}
  {\haalbuffer}
\voorbeeld[vb:zesring]
\typebuffer
\stopfiguurtekst

Met \type{\stelchemiein} kunnen verschillende kenmerken van het zetwerk
worden ingesteld. De met dit commando ingestelde waarden gelden voor alle
volgende formules.\voetnoot{Natuurlijk kan de scope worden beperkt met behulp
van \type{{ }} en de groeperingsmacro's \type{..group}. De instellingen
kunnen ook direct achter \tex{startchemie} worden opgegeven. In dat geval
blijven de instellingen beperkt tot ��n structuurformule.} 

\startbuffer
\startchemie[kader=aan,breedte=passend,hoogte=passend]
  \chemie[CARBON,CB1][A,B,C,D,E,F]
\stopchemie
\stopbuffer

\startfiguurtekst
  [links][]
  {geen}
  {\haalbuffer}
\voorbeeld
\typebuffer
\stopfiguurtekst

Zoals uit beide voorbeelden blijkt, is \type{\chemie} het centrale commando.
Dit commando, dat meerdere malen binnen een \type{\start}||\type{\stop}||paar
kan worden opgegeven, krijgt een of twee argumenten mee. Deze worden
tussen~\type{[ ]} opgegeven. Het eerste argument heeft betrekking op de te
tekenen bindingen, het tweede bevat de weer te geven atomen of moleculen.
Tekst wordt in de wiskundige mode gezet, dat wil zeggen dat alles wat normaal
gesproken tussen~\type{$ $} is toegestaan, mag worden opgegeven. 

We werken hier het eerste voorbeeld uit. Allereerst is het trefwoord
\type{SIX} meegegeven. Hiermee geven we aan dat we een zesring tekenen.
Analoog kennen we \type{ONE}, \type{THREE}, \type{FOUR} en \type{FIVE},
\type{EIGHT}, \type{CARBON}, \type{NEWMAN}, \type{CHAIR}, enkele varianten
hierop en wat symbolen. 

\def\structuurtje[#1]%
  {\vbox
     {\startchemie[kader=aan,schaal=klein,formaat=klein]%
      \chemie[#1][1,2,3,4,5,6,7,8]%
      \stopchemie}}

\plaatsfiguur{geen}
 \startcombinatie[4*3]
   {\structuurtje[ONE,SB,Z]}           {One}
   {\structuurtje[THREE,B,R,RZ]}       {Three}
   {\structuurtje[FOUR,B,R,RZ]}        {Four}
   {\structuurtje[FIVE,B,R,CRZ]}       {Five}
   {\structuurtje[SIX,B,R,RZ]}         {Six}
   {\structuurtje[EIGHT,B]}            {Eight}
   {\structuurtje[FIVE,FRONT,B,-R,+R]} {Five Front}
   {\structuurtje[SIX,FRONT,B,-R,+R]}  {Six Front}
   {\structuurtje[CARBON,CB]}          {Carbon}
   {\structuurtje[NEWMAN,STAGGER,CB]}  {Newman Stacker}
   {\structuurtje[NEWMAN,ECLIPSE,CB]}  {Newman Eclipse}
   {\structuurtje[CHAIR,B,-R,+R]}      {Chair}
 \stopcombinatie

Qua afmetingen spring \type{CHAIR} er duidelijk uit. Ook in andere opzichten
wijkt deze structuur wat af van de andere structuren. Zo is hier roteren en
combineren niet mogelijk. 

\startbuffer
\startchemie[kader=aan,schaal=klein]
  \chemie
    [CHAIR,B,+R,-R,
     +RZ1,+RZ2,+RZ3,+RZ4,+RZ5,+RZ6,
     -RZ1,-RZ2,-RZ3,-RZ4,-RZ5,-RZ6]
    [a,b,c,d,e,f,1,2,3,4,5,6]
\stopchemie
\stopbuffer

\startfiguurtekst
  [links][]
  {geen}
  {\haalbuffer}
\voorbeeld
\typebuffer
\stopfiguurtekst

Binnen structuren worden de chemische bindingen tussen de \kap{c}||atomen op
een vergelijkbare wijze aangegeven. Zo gebruiken we in het voorbeeld~\type{B}
en~\type{R}. Bindingen zijn genummerd en kunnen op verschillende manieren
worden opgegeven: 

\starttypen
\chemie[SIX,B1,B2,B3,B4,B5,B6]
\chemie[SIX,B135]
\chemie[SIX,B1..5]
\stoptypen

Deze commando's tekenen delen van een zesring. Met \type{R} en \type{RZ}
kunnen we substituenten aan de zesring plaatsen. Het commando \type{R} tekent
vanuit een hoekpunt van de zesring de aanzet tot een binding met een
substituent ($\angle~120^\circ$). Het betreffende hoekpunt wordt aangegeven
met een getal. 

\starttypen
\chemie[SIX,B1..6,R1..6]
\stoptypen

De bovenstaande aanroep plaatst alleen de binding naar de substituenten. De
substituent zelf wordt met \type{RZ} aangegeven. Ook hier worden getallen
gebruikt om de positie te markeren. De substituenten worden in het tweede,
optionele argument als tekst meegegeven. 

\startbuffer
\startchemie[kader=aan,breedte=6000]
  \chemie[SIX,B1..6,R1..6,RZ1..3][CH_3,CH_3,OH]
\stopchemie
\stopbuffer

\startfiguurtekst
  [links][]
  {geen}
  {\haalbuffer}
\voorbeeld
\typebuffer
\stopfiguurtekst

Als het tweede argument wordt weggelaten, worden geen teksten geplaatst en
heeft het commando \type{RZ1..3} geen gevolg. 

\hoofdstuk{Bindingen}

Hieronder zijn overzichten opgenomen van de bindingen die men kan aantreffen
bij de verschillende structuren. Uit de voorbeelden en overzichten verderop
in deze handleiding zal blijken waar de commando's voor staan. 

In de linker kolom staan steeds de volledige bindingen weergegeven, in de
rechter kolom de ingekorte bindingen. Deze laatste maken het mogelijk atomen
en moleculen in de binding op te nemen. Een aantal bindingen kan aan beiden
kanten, links (\type{-}) of rechts (\type{+}) worden ingekort. 

\plaatstabel
{Enkelvoudige bindingen.}
\starttabel[|rT|lw(4cm)|rT|lw(4cm)|]
\HL
\VL ~~~B \VL Bond          \VL ~~SB \VL Single Bond       \VL\FR
\VL ~~BB \VL Bold Bond     \VL ~-SB \VL Left Single Bond  \VL\MR
\VL ~~HB \VL Hydrogen Bond \VL ~+SB \VL Right Single Bond \VL\LR
\HL
\stoptabel

In het onderstaande voorbeeld zien we een aantal bindingen
gecombineerd:

\startbuffer
\startchemie[kader=aan]
  \chemie[ONE,SD1,SB4,BB2,SB7,Z01247][C,H,H,H,H]
\stopchemie
\stopbuffer

\startfiguurtekst
  [links][]
  {geen}
  {\haalbuffer}
\voorbeeld
\typebuffer
\stopfiguurtekst

Een binding kan worden gevolgd door een of meer getallen of een range,
bijvoorbeeld: \type{B1}, \type{B135} en \type{B1..5}. Als alle bindingen
nodig zijn, kan worden volstaan met~\type{B}. 

Binnen een ring kan een extra binding worden aangegeven en tussen atomen en
moleculen dubbele of drievoudige bindingen. 

\plaatstabel
{Meervoudige bindingen.}
\starttabel[|rT|lw(4cm)|rT|lw(4cm)|]
\HL
\VL ~~EB \VL Extra Bond    \VL ~~DB \VL Double Bond       \VL\FR
\VL      \VL               \VL ~~TB \VL Triple Bond       \VL\LR
\HL
\stoptabel

We kunnen extra ladingen op verschillende manieren oproepen. De betreffende
commando's beginnen met een \type{E}. 

\plaatstabel
{Ladingen.}
\starttabel[|rT|lw(4cm)|rT|lw(4cm)|]
\HL
\VL ~~ES \VL Extra Single  \VL ~~ED \VL Extra Double \VL\FR
\VL ~~EP \VL Extra Pair    \VL ~~ET \VL Extra Triple \VL\LR
\HL
\stoptabel

Het onderstaande koolstofatoom is uit demonstratieoogpunt voorzien van 
maar liefst 8~buitenelektronen. 

\startbuffer
\startchemie[kader=aan]
  \chemie[ONE,Z0,ES1,ED3,ET5,EP7][C]
\stopchemie
\stopbuffer

\startfiguurtekst
  [links][]
  {geen}
  {\haalbuffer}
\voorbeeld
\typebuffer
\stopfiguurtekst

Een binding kan worden kortgesloten. Dit komt voor bij bijvoorbeeld
zesringen. In dat geval wordt het atoom dat moet worden overgeslagen
opgegeven. Daarnaast kan binnen een zesring een cirkel worden getekend. 

\plaatstabel
{Bijzondere bindingen.}
\starttabel[|rT|lw(4cm)|rT|lw(4cm)|]
\HL
\VL ~~SS \VL Short Shortcut       \VL ~~~S \VL Shortcut            \VL\FR
\VL ~-SS \VL Left Short Shortcut  \VL ~MID \VL Open Mid Shortcut   \VL\MR
\VL ~+SS \VL Right Short Shortcut \VL MIDS \VL Closed Mid Shortcut \VL\LR
\HL
\stoptabel

\plaatstabel
{Circelvormige bindingen.}
\starttabel[|rT|lw(4cm)|rT|lw(4cm)|]
\HL
\VL ~~~C \VL Circle               \VL ~~CD \VL Dashed Circle          \VL\FR
\VL ~~CC \VL Shifted Circle       \VL ~CCD \VL Dashed Shifted Circle  \VL\LR
\HL
\stoptabel

Hoewel de circelvormige \type{C..} commando's zich laten raden, tonen we hier
toch een voorbeeld.

\startbuffer
\startchemie[kader=aan]
  \chemie
    [SIX,B,ER6,CCD12346,Z0,PB:RZ6,ONE,SB8,EP6,Z0,ZT6,Z8,PE]
    [\ominus,O,\oplus,H]
\stopchemie
\stopbuffer

\startfiguurtekst
  [links][]
  {geen}
  {\haalbuffer}
\voorbeeld
\typebuffer
\stopfiguurtekst

Aan alle hoekpunten kunnen substituenten worden verbonden. Het begrip
substituent mag hier overigens ruim worden opgevat. Afhankelijk van de
aanwezigheid van atomen en moleculen, kunnen de bindingen kort of lang zijn.
We zullen deze commando's vaak tegenkomen. 

\plaatstabel
{Bindingen naar substituenten.}
\starttabel[|rT|lw(4cm)|rT|lw(4cm)|]
\HL
\VL ~~~R \VL Radical       \VL ~~SR \VL Single Radical       \VL\FR
\VL ~~-R \VL Left Radical  \VL ~-SR \VL Single Left Radical  \VL\MR
\VL ~~+R \VL Right Radical \VL ~+SR \VL Single Right Radical \VL\LR
\HL
\stoptabel

Naast deze eenvoudige lijntjes zijn er wat alternatieven om bruggen aan te
geven. 

\plaatstabel
{Bijzondere substituenten.}
\starttabel[|rT|lw(4cm)|rT|lw(4cm)|]
\HL
\VL ~~RD \VL Radical Dashed       \VL ~~RB\VL Radical Bold       \VL\FR
\VL ~-RD \VL Left Radical Dashed  \VL ~-RB\VL Left Radical Bold  \VL\MR
\VL ~+RD \VL Right Radical Dashed \VL ~+RB\VL Right Radical Bold \VL\LR
\HL
\stoptabel

Radicalen kunnen steeds op drie manieren worden getekend. De linker en
rechter variant zullen echter zelden worden gebruikt. 

\startbuffer
\startchemie[kader=aan]
  \chemie[SIX,B,R14,RD25,RB36]
\stopchemie
\stopbuffer

\startfiguurtekst
  [links][]
  {geen}
  {\haalbuffer}
\voorbeeld
\typebuffer
\stopfiguurtekst

\plaatstabel
{Nog meer bijzondere substituenten.}
\starttabel[|rT|lw(4cm)|rT|lw(4cm)|]
\HL
\VL ~~SD \VL Single Dashed \VL ~LDD \VL Left Double Dashed  \VL\FR
\VL ~~OE \VL Open Ended    \VL ~RDD \VL Right Double Dashed \VL\LR
\HL
\stoptabel

Hieronder is een voorbeeld van een {\em Open End} weergegeven. We zien hier
een zesring (\type{SIX}) met daaraan een aaneenschakeling van \type{ONE}'s.
Lees voor het moment maar even over het gebruik van \type{PB} heen. 

\startbuffer
\startchemie
    [breedte=5000,boven=2500,onder=1500,kader=aan]
  \chemie
    [SIX,B,C,R6,
     PB:RZ6,ONE,CZ0,OE1,SB5,MOV5,CZ0,OFF5,OE5,PE]
    [CH,CH_2]
\stopchemie
\stopbuffer

\startfiguurtekst
  [links][]
  {geen}
  {\haalbuffer}
\voorbeeld
\typebuffer
\stopfiguurtekst

Natuurlijk kunnen substituenten ook door dubbele bindingen aan de structuur
worden verbonden. 

\plaatstabel
{Dubbele bindingen naar substituenten.}
\starttabel[|rT|lw(4cm)|rT|lw(4cm)|]
\HL
\VL ~~ER \VL Extra Radical \VL ~~DR \VL Double Radical \VL\SR
\HL
\stoptabel

Aan bindingen kunnen teksten worden gekoppeld. Deze teksten worden in de
opgegeven volgorde uit de tweede set achter \type{\chemie} gehaald. 

\plaatstabel
{Atomen en moleculen (radikalen).}
\starttabel[|rT|lw(4cm)|rT|lw(4cm)|]
\HL
\VL ~~~Z \VL Atom        \VL ~~RZ \VL Radical Atom       \VL\FR
\VL ~CRZ \VL Center Atom \VL ~-RZ \VL Left Radical Atom  \VL\MR
\VL MIDZ \VL Mid Atom    \VL ~+RZ \VL Right Radical Atom \VL\LR
\HL
\stoptabel

Van deze commando's is \type{RZ}, dat aansluit op \type{R}, de meest voor de
hand liggende. Het eerder genoemde \type{MID} commando is alleen beschikbaar
bij een zesring (\type{SIX}). In het voorbeeld hieronder zien we zowel
\type{MID} als \type{MIDZ} in actie. Deze commando's krijgen geen positie
mee. 

\startbuffer
\startchemie[kader=aan]
  \chemie[SIX,B,MID,MIDZ][\SL{CH_2}]
\stopchemie
\stopbuffer

\startfiguurtekst
  [links][]
  {geen}
  {\haalbuffer}
\voorbeeld
\typebuffer
\stopfiguurtekst

De atomen/moleculen worden met de klok mee genummerd. Ook hier mogen
combinaties worden opgegeven. De positie \type{0} (nul) valt samen met het
midden van een structuur. 

We kunnen aan een atoom of binding labels en nummers toevoegen. Dit doen we
met \type{ZN} en \type{ZT}: 

\plaatstabel
{Labels en nummers.}
\starttabel[|rT|lw(4cm)|rT|lw(4cm)|]
\HL
\VL ~~ZN \VL Atom Number \VL ~~ZT \VL Atom Text \VL\SR
\HL
\stoptabel

Bij \type{SIX} en \type{FIVE} kan dat ook bij radikalen. In dat geval
gebruiken we \type{RN} en \type{RT}. 

\plaatstabel
{Labels en nummers.}
\starttabel[|rT|lw(4cm)|rT|lw(4cm)|]
\HL
\VL ~~RN \VL Radical Number        \VL ~~RT \VL Radical Text        \VL\FR
\VL ~RTN \VL Radical Top Number    \VL ~RTT \VL Radical Top Text    \VL\MR
\VL ~RBN \VL Radical Bottom Number \VL ~RBT \VL Radical Bottom Text \VL\LR
\HL
\stoptabel

In geval van \type{ONE} hebben we bovendien nog boven (top) en onder (bottom)
varianten. 

\plaatstabel
{Extra labels en nummers.}
\starttabel[|rT|lw(4cm)|rT|lw(4cm)|]
\HL
\VL ~ZTN \VL Atom Top Number    \VL ~ZTT \VL Atom Top Text    \VL\FR
\VL ~ZBN \VL Atom Bottom Number \VL ~ZBT \VL Atom Bottom Text \VL\LR
\HL
\stoptabel

Bij \type{ZTN} en \type{ZBN} worden de nummers automatisch gegenereerd. De
andere commando's gebruiken de opgegeven tekst 

\startbuffer
\startchemie[kader=aan]
  \chemie[ONE,SB,Z0,ZTT][C,a,b,c,d,e,f,g,h]
\stopchemie
\stopbuffer

\startfiguurtekst
  [links][]
  {geen}
  {\haalbuffer}
\voorbeeld
\typebuffer
\stopfiguurtekst

Resten ons nog wat symbolen om de structuur aan te kleden.

\plaatstabel
{Indicaties.}
\starttabel[|rT|lw(4cm)|rT|lw(4cm)|]
\HL
\VL ~~AU \VL Arrow Up \VL ~~AD \VL Arrow Down \VL\SR
\HL
\stoptabel

De pijlen worden getekend tussen de radicalen.

\startbuffer
\startchemie[kader=aan]
  \chemie[SIX,B,AU]
\stopchemie
\stopbuffer

\startfiguurtekst
  [links][]
  {geen}
  {\haalbuffer}
\voorbeeld
\typebuffer
\stopfiguurtekst

We merken nog op dat bij het plaatsen van atomen en moleculen zo goed
mogelijk wordt rekening gehouden met de (mogelijke) afmetingen van atomen en
moleculen. De breedte van de \chemie{C} en de hoogte van \chemie{C^n_m}
spelen daarbij een rol. Dit mechanisme kan nog worden verfijnd. 

\hoofdstuk{Vooraanzichten}

De structuren \type{FIVE} en \type{SIX} zijn ook als vooraanzicht
beschikbaar, zij het met wat beperkingen. We tonen hieronder twee voorbeelden
van zo'n vooraanzicht. 

Vooraanzichten kunnen niet worden geroteerd. Ook zijn er beperkingen ten
aanzien van het koppelen aan andere structuren. Dit hangt samen met het
karakter van de weergave. 

\startbuffer
\startchemie[hoogte=4500,onder=2500]
  \bottext{$\beta$-D-Fructopyranose}
  \chemie
    [SIX,FRONT,BB1236,+SB4,-SB5,Z5][O]
  \chemie
    [SIX,FRONT,+R12346,+RZ12346][\SR{HO},H,H,H,OH]
  \chemie
    [SIX,FRONT,-R12346,-RZ12346][H,OH,\SR{HO},H,CH_2OH]
\stopchemie
\stopbuffer

\startfiguurtekst
  [links][]
  {geen}
  {\haalbuffer}
\voorbeeld
\typebuffer
\stopfiguurtekst

De positionering van de radicalen is een mix tussen haalbaarheid en
kwaliteit. In het volgende voorbeeld blijkt dit nog duidelijker. 

\startbuffer
\startchemie[hoogte=4500,onder=2500]
  \bottext{$\alpha$-D-Fructofuranose}
  \chemie
    [FIVE,FRONT,BB125,+SB3,-SB4,Z4][O]
  \chemie
    [FIVE,FRONT,+R1235,+RZ1235][\SR{HO},H,\SR{HOH_2C},CH_2OH]
  \chemie
    [FIVE,FRONT,-R1235,-RZ1235][OH,H,\SR{HO},H,CH_2OH]
\stopchemie
\stopbuffer

\startfiguurtekst
  [links][]
  {geen}
  {\haalbuffer}
\voorbeeld
\typebuffer
\stopfiguurtekst

\hoofdstuk{Definities} 

Het is mogelijk een bibliotheek van structuren op te bouwen. Deze structuren
kunnen we, al naar gelang de behoefte, later oproepen en voorzien van extra
componenten. Ook kunnen ze dienen als bouwstenen voor ingewikkelder
structuren. Het voordefini�ren van structuren kan plaatsvinden met behulp van
de \TEX||primitieve \type{\def}. 

\starttypen
\def\zesring{\chemie[SIX,B,R,RZ]}
\stoptypen

In plaats van \type{\def} kan ook het onderstaande commando worden gebruikt.
In dat geval wordt een melding gegeven als de definitie reeds bestaat. 

\starttypen
\definieerchemie[zesring]
  {\chemie[SIX,B,R,RZ]}
\stoptypen

Hoewel beide manieren van defini�ren zijn toegestaan is de tweede manier
robuuster, omdat in dat geval beschermende maatregelen worden genomen om
conflicten met bestaande commando's te voorkomen. 

De aanroep \type{\chemie[zesring]} levert een zesring op zonder
substituenten. Er is immers geen tweede argument gegeven. 

\startbuffer
\definieerchemie[zesring]
  {\chemie[SIX,B,R,RZ]}

\startchemie[kader=aan,breedte=6000]
  \chemie[zesring]
\stopchemie
\stopbuffer

\startfiguurtekst
  [links][]
  {geen}
  {\haalbuffer}
\voorbeeld
\typebuffer
\stopfiguurtekst

Als we zes substituenten willen plaatsen, dan kan dat als volgt: 

\startbuffer
\definieerchemie[zesring]
  {\chemie[SIX,B,R,RZ]}

\startchemie[kader=aan,breedte=6000]
  \chemie[zesring][R_1,R_2,R_3,R_4,R_5,R_6]
\stopchemie
\stopbuffer

\startfiguurtekst
  [links][]
  {geen}
  {\haalbuffer}
\voorbeeld
\typebuffer
\stopfiguurtekst

De structuur \type{zesring} kan ook zonder substituenten (\type{RZ}) worden
gedefinieerd. In dat geval worden bij de aanroep \type{\chemie[zesring]} geen
substituenten verwacht. Dat we deze alsnog kunnen plaatsen toont het
onderstaande voorbeeld. 

\startbuffer
\definieerchemie[zesring]
  {\chemie[SIX,B,R]}

\startchemie[kader=aan,breedte=6000]
  \chemie[zesring,RZ][A,B,C,D,E,F]
\stopchemie
\stopbuffer

\startfiguurtekst
  [links][]
  {geen}
  {\haalbuffer}
\voorbeeld
\typebuffer
\stopfiguurtekst

In principe is het aantal mogelijkheden onbegrensd. Men dient zich echter
steeds te realiseren dat de atomen en moleculen uit het tweede argument
worden opgehaald in de volgorde van het eerste argument. 

In een definitie mogen ook atomen en moleculen (teksten) worden geplaatst. 

\starttypen
\definieerchemie[zesring]
  {\chemie[SIX,B,R,RZ135][R_1,R_3,R_5]}
\stoptypen

Hier worden dus altijd drie substituenten geplaatst. Als we bij het oproepen
meer substituenten willen plaatsen, dan dienen we expliciet aan te geven dat
we doorgaan met de zesring (\type{SIX}). 

\startbuffer
\definieerchemie[zesring]
  {\chemie[SIX,B,R,RZ135][R_1,R_3,R_5]}

\startchemie[kader=aan,breedte=6000]
  \chemie[zesring,SIX,RZ246][A,B,C]
\stopchemie
\stopbuffer

\startfiguurtekst
  [links][]
  {geen}
  {\haalbuffer}
\voorbeeld
\typebuffer
\stopfiguurtekst

In definities werkt \type{\chemie[]} dus globaal (dat wil zeggen dat
\type{SIX} wordt onthouden) en \type{\chemie[][]} lokaal. De idee hierachter
is dat in het eerste geval een reeks commando's wordt tussengevoegd en in het
tweede geval een complete, zelfstandige structuur. 

In een definitie kan dus meerdere malen \type{\chemie} voorkomen. Het vorige
voorbeeld had daarom ook kunnen worden opgeroepen met: 

\startbuffer
\definieerchemie[zesring]
  {\chemie[SIX,B,R,RZ135][R_1,R_3,R_5]
   \chemie[SIX,RZ246]}

\startchemie[kader=aan,breedte=6000]
  \chemie[zesring][A,B,C]
\stopchemie
\stopbuffer

\startfiguurtekst
  [links][]
  {geen}
  {\haalbuffer}
\voorbeeld
\typebuffer
\stopfiguurtekst

Als \TEX\ met de melding komt dat er sprake is van een \type{onbekend
commando}, dan is men waarschijnlijk vergeten \type{SIX}, \type{FIVE} of een
vergelijkbaar structuur||commando mee te geven. 

\hoofdstuk{Combinaties}

Structuren kunnen worden gecombineerd tot complexe verbindingen. Het
verplaatsen van de ene structuur ten opzichte van de andere structuur gebeurt
met \type{MOV}, \type{ROT}, \type{ADJ} en \type{SUB}. 

\plaatstabel
{Verplaatsingen en rotaties.}
\starttabel[ | r T | l w(2cm) | p(8cm) | ]
\HL
\VL MOV \VL Move       \VL het verplaatsen van een zelfde \hbox{structuur} in
                           de richting van een binding\VL\FR 
\VL ADJ \VL Adjace     \VL het verplaatsen van een andere structuur in de
                           richting van de $x$-- of $y$--as, aanliggend aan
                           een binding\VL\MR 
\VL SUB \VL Substitute \VL het verplaatsen van de ene structuur ten opzichte
                           van een andere in de richting van de $x$-- of
                           $y$--as\VL\MR 
\VL ROT \VL Rotate     \VL het roteren van een structuur\VL\LR 
\HL
\stoptabel

De bovenstaande vier commando's hebben binnen de verschillende structuren een
ander effect. Zo is de hoek waarover wordt geroteerd bij
\type{\chemie[FIVE,ROT1,B]} anders dan die bij \type{\chemie[SIX,ROT1,B]}. 

Bij \type{ONE} is naast \type{MOV} ook \type{DIR} beschikbaar. Beide
commando's hebben dezelfde werking maar verschillen in spatiering. Kleine
verplaatsingen zijn daarnaast mogelijk, met \type{OFF}. 

\plaatstabel
{Verplaatsingen en rotaties.}
\starttabel[ | r T | l w(2cm) | p(8cm) | ]
\HL
\VL DIR \VL Direction \VL Een \type{MOV}, bedoeld voor diagonaal positioneren 
                          van structuren\VL\FR
\VL OFF \VL Offset    \VL Een kleine verplaatsing van een atoom of molecuul\VL\LR 
\HL
\stoptabel

Binnen \type{CARBON} is het bovendien mogelijk een structuur te spiegelen.
Dit gebeurt met \type{MIR}. 

\plaatstabel
{Spiegelen.}
\starttabel[ | r T | l w(2cm) | p(8cm) | ]
\HL
\VL MIR \VL Mirror     \VL het spiegelen van de structuur\VL\SR
\HL
\stoptabel

Met een cijfer geven we de richting van een verplaatsing of de mate van de
rotatie aan. Omdat deze commando's nauw verbonden zijn met de actuele
structuur, dienen ze te worden gegeven voordat bindingen en teksten worden
getekend. Het maakt dus uit of \type{\chemie[FIVE,B,ROT1,R]} wordt opgegeven
of \type{\chemie[FIVE,ROT1,B,R]}. De eerste aanroep levert een ongewenst
resultaat. 

\startbuffer
\startchemie[kader=aan,breedte=4000,rechts=3000]
  \chemie[SIX,B,MOV1,B]
\stopchemie
\stopbuffer

\startfiguurtekst
  [links][]
  {geen}
  {\haalbuffer}
\voorbeeld
\typebuffer
\stopfiguurtekst

Achtereenvolgens wordt hier een zesring getekend: \type{SIX,B}, een
verplaatsing in de richting van binding~1 gerealiseerd: \type{MOV1} en een
tweede zesring getekend:~\type{B}. Een verplaatsing met \type{MOV} kan bij
een zesring in zes richtingen plaatsvinden. Dit in tegenstelling tot een
verplaatsing met \type{ADJ}, die in de vier as||richtingen plaatsvindt ($x$,
$-x$, $y$, $-y$). Bij een zesring vallen enkele van deze verplaatsingen
samen: het bovenstaande voorbeeld had ook kunnen worden gerealiseerd met:
\type{[SIX,B,ADJ1,B]}. 

Ook kunnen verschillende structuren worden gecombineerd. Aan een structuur
\type{FIVE} kan bijvoorbeeld \type{SIX} worden gekoppeld. Het mechanisme dat
voor de koppeling zorgt is voor de gebruiker grotendeels verborgen. In het
volgende voorbeeld wordt achtereenvolgens een zesring getekend: \type{SIX,B},
een verplaatsing langs de positieve $x$||as gerealiseerd: \type{ADJ1}, en een
geroteerde vijfring getekend: \type{FIVE,ROT3,B}. 

\startbuffer
\startchemie[kader=aan,breedte=4000,rechts=3000]
  \chemie[SIX,B,ADJ1,FIVE,ROT3,B]
\stopchemie
\stopbuffer

\startfiguurtekst
  [links][]
  {geen}
  {\haalbuffer}
\voorbeeld
\typebuffer
\stopfiguurtekst

\startbuffer
\startchemie[kader=aan,breedte=5000,rechts=4500]
  \chemie[SIX,ROT2,B,R6,SUB1,FIVE,B,R4]
\stopchemie
\stopbuffer

\startfiguurtekst
  [links][]
  {geen}
  {\haalbuffer}
\voorbeeld
\typebuffer
\stopfiguurtekst

Een overgang naar een aansluitende structuur vindt dus plaats met \type{ADJ}.
Vaak zal, om een goede aansluiting te krijgen, een van de twee structuren
moeten worden geroteerd met \type{ROT}. Als een structuur niet direct maar
via een binding wordt gekoppeld, gebruikt men \type{SUB}. Rotaties vinden
plaats in stappen van 90\graden, met de klok mee. Verplaatsingen met
\type{ADJ} en \type{SUB} vinden plaats in de vier asrichtingen. 

De volgende voorbeelden illustreren nog eens duidelijk dat de afmetingen van
de kleinere structuren bepaald zijn door die van de grote, met name
\type{SIX}. De hieronder getoonde \type{EIGHT} heeft overigens wat minder
mogelijkheden heeft dan bijvoorbeeld \type{SIX}. 

\startbuffer
\startchemie[breedte=6000,links=1500,kader=aan]
  \chemie[EIGHT,B,MOV1,B]
\stopchemie
\stopbuffer

\startfiguurtekst
  [links][]
  {geen}
  {\haalbuffer}
\voorbeeld
\typebuffer
\stopfiguurtekst

\startbuffer
\startchemie[breedte=6000,links=1500,kader=aan]
  \chemie[EIGHT,B,ADJ1,SIX,B]
\stopchemie
\stopbuffer

\startfiguurtekst
  [links][]
  {geen}
  {\haalbuffer}
\voorbeeld
\typebuffer
\stopfiguurtekst

\startbuffer
\startchemie[breedte=6000,links=1500,kader=aan]
  \chemie[EIGHT,B,ADJ1,FIVE,ROT3,B]
\stopchemie
\stopbuffer

\startfiguurtekst
  [links][]
  {geen}
  {\haalbuffer}
\voorbeeld
\typebuffer
\stopfiguurtekst

Inmiddels zal duidelijk zijn dat het uitmaakt in welke volgorde we de
commando's geven. Een voor de hand liggende volgorde is: 

\starttypen
\chemie
  [structuur,                        % SIX, FIVE, ...
   bindingen binnen de structuur,    % B, C, EB, ...
   bindingen buiten de structuur,    % R, DR, ...
   te plaatsen atomen,               % Z
   te plaatsen subsituenten]         % RZ, -RZ, ...
  [atomen,
   substituenten]
\stoptypen

Het aaneenschakelen van structuren komt in de regel neer op transleren en
roteren. Hoewel dit misschien niet direct zo lijkt, zit hierin een zekere
systematiek. Het proces zou dan ook kunnen worden vereenvoudigd. De in
experimentele versies reeds gerealiseerde automatisering is weer ongedaan
gemaakt, omdat gebleken is dat 'verborgen' rotaties leiden tot misverstanden
met betrekking tot de plaats van bindingen. Bovendien is het eenvoudiger een
niet geroteerde structuur van bindingen, atomen en moleculen te voorzien dan
een geroteerde. Een samengestelde structuur kan beter eerst per onderdeel
worden gedefinieerd, eventueel met translaties, en pas als laatste stap
worden geroteerd. 

Vaak zal een zesring een of meerdere uitlopers samengesteld uit \type{ONE}
hebben. Dergelijke ketens zig||zaggen als het ware om een centrale as. In
zulke situaties gebruiken we \type{DIR}. 

\startbuffer
\startchemie
  [schaal=klein,breedte=6000,hoogte=6000,kader=aan]
  \chemie[SIX,SB2356,DB14,Z2346,SR36,RZ36]
    [C,N,C,C,H,H_2]
  \chemie[PB:Z1,ONE,Z0,DIR8,Z0,SB24,DB7,Z27,PE]
    [C,C,CH_3,O]
  \chemie[PB:Z5,ONE,Z0,DIR6,Z0,SB24,DB7,Z47,PE]
    [C,C,H_3C,O]
  \chemie[SR24,RZ24]
    [CH_3,H_3C]
\stopchemie
\stopbuffer

\startfiguurtekst
  [links][]
  {geen}
  {\haalbuffer}
\voorbeeld
\typebuffer
\stopfiguurtekst

Omdat de ketens geen vaststaande formaten hebben, worden ze opgebouwd en
gepositioneerd als substructuur. We gebruiken voor de positionering de
commando's \type{PB} en \type{PE}. 

\plaatstabel
{Positioneren.}
\starttabel[ | r T | l w(4cm) | p(6cm) | ]
\HL
\VL PB:.. \VL Picture Begin \VL het starten van een substructuur   \VL\FR
\VL PE    \VL Picture End   \VL het afsluiten van een substructuur \VL\LR
\HL
\stoptabel

Direct na \type{PB} volgt de positie waarop de substructuur wordt geplaatst.
Het is belangrijk in de gaten te houden dat het eerstvolgende atoom
gecentreerd wordt op deze positie. Gebruik daarom bij voorkeur een centraal
atoom. 

De aanleiding voor het introduceren van deze commando's lag in het feit dat
de getoonde structuur moet kunnen worden gezet op een optisch acceptabele
wijze. Er zijn echter verschillende manieren om dergelijke structuren te
maken. De volgende variant levert ook resultaat: 

\startbuffer
\startchemie
    [schaal=klein,breedte=6000,hoogte=6000,kader=aan]
  \chemie
    [SIX,SB2356,DB14,Z36,SR36,RZ36][N,C,H,H_2]
  \chemie
    [PB:Z1,ONE,Z0,DIR8,Z0,SB24,DB7,Z27,PE][C,C,CH_3,O]
  \chemie
    [PB:Z5,ONE,Z0,DIR6,Z0,SB24,DB7,Z47,PE][C,C,H_3C,O]
  \chemie
    [PB:Z2,ONE,Z0,DIR2,SB6,CZ0,PE][C,CH_3]
  \chemie
    [PB:Z4,ONE,Z0,DIR4,SB8,CZ0,PE][C,H_3C]
\stopchemie
\stopbuffer

\startfiguurtekst
  [links][]
  {geen}
  {\haalbuffer}
\voorbeeld
\typebuffer
\stopfiguurtekst

De meest effici�nte wijze om een dergelijke structuur te defini�ren is echter
de volgende. Deze laatste oplossing is niet bij voorbaat de typografisch
mooiste. 

\startbuffer
\startchemie
    [schaal=klein,breedte=6000,hoogte=6000,kader=aan]
  \chemie[SIX,SB2356,DB14,Z,SR36,RZ36,SR1245,RZ24]
    [C,C,N,C,C,C,H,H_2,CH_3,H_3C]
  \chemie[PB:RZ1,ONE,Z0,SB2,DB7,Z27,PE]
    [C,CH_3,O]
  \chemie[PB:RZ5,ONE,Z0,SB4,DB7,Z47,PE]
    [C,H_3C,O]
\stopchemie
\stopbuffer

\startfiguurtekst
  [links][]
  {geen}
  {\haalbuffer}
\voorbeeld
\typebuffer
\stopfiguurtekst

Het is misschien al opgevallen dat de maat van de structuur wordt bepaald
door de substituenten. De ketens tellen, afgezien van het start||atoom, niet
mee. Dit bevordert een consistente opbouw. 

De verschillen tussen \type{SUB} en \type{PB} zijn subtiel maar niet te
verwaarlozen. Vergelijk bijvoorbeeld de volgende formules: 

\startbuffer
\startchemie
  [breedte=passend,kader=aan,schaal=klein]
  \chemie
    [SIX,ROT2,B,C,R236,RZ23,
     SUB1,ONE,OFF1,Z0,4OFF1,SB1,Z1]
    [HO,HO,CHOH,CH_2NH_2]
\stopchemie
\stopbuffer

\startfiguurtekst
  [links][]
  {geen}
  {\haalbuffer}
\voorbeeld
\typebuffer
\stopfiguurtekst

\startbuffer
\startchemie
  [breedte=passend,kader=aan,schaal=klein]
  \chemie
    [SIX,ROT2,B,C,R236,RZ23,
     PB:RZ6,ONE,Z0,3OFF1,SB1,Z1,PE]
    [HO,HO,CHOH,CH_2NH_2]
\stopchemie
\stopbuffer

\startfiguurtekst
  [links][]
  {geen}
  {\haalbuffer}
\voorbeeld
\typebuffer
\stopfiguurtekst

Het gebruik van \type{PB:} is wat lastiger, maar geeft mooiere resultaten. We
zien hier overigens duidelijk dat we in dat geval de breedte zelf moeten
opgeven. Dit komt omdat substituenten nooit meetellen bij de beoaling van de
afmetingen. Zou dit wel het geval zijn, dan zouden in veel gevallen
structuren inconsistente afmetingen krijgen. Ze zijn dan moeilijk te
combineren. 

We gaan eerst wat neder in op \type{OFF}. Bij \type{ONE} kan \type{Z0} uit
meer dan een atoom bestaan. In dat geval is de gereserveerde ruimte
ontoereikend. Als voor \type{Z0} meer ruimte nodig is, dan kunnen de
bindingen~1, 2 en~8 worden opgeschoven met het commando \type{OFF}, wat staat
voor 'offset'. Hoewel we dit commando al in eerdere voorbeelden zagen, is
hieronder nog een voorbeeld opgenomen. 

\startbuffer
\startchemie
  [breedte=passend,formaat=klein,schaal=klein,kader=aan]
  \chemie
    [SIX,B,C,ADJ1,
     FIVE,ROT3,SB34,+SB2,-SB5,Z345,DR35,SR4,CRZ35,SUB1,
     ONE,OFF1,SB258,Z0,Z28]
    [C,N,C,O,O,
     CH,COOC_2H_5,COOC_2H_5]
\stopchemie
\stopbuffer

\startfiguurtekst
  [links][]
  {geen}
  {\haalbuffer}
\voorbeeld
\typebuffer
\stopfiguurtekst

Een verplaatsing biedt ruimte voor een extra karakter. We hadden meer ruimte
gekregen als we bijvoorbeeld \type{3OFF1} hadden gebruikt. Dergelijke, op het
eerste oog vrij ingewikkelde, definities kunnen worden gerealiseerd door
eerst de afzonderlijke delen te defini�ren. Het roteren kan daarbij voor het
laatst worden bewaard. 

Er duikt hier nog een nieuw commando op: \type{CRZ}. Dit commando kan worden
gebruikt om een atoom of molecuul in het verlengde van de binding te
plaatsen, wat in dit geval gewenst is. Hetzelfde had kunnen worden bereikt
met het commando \type{RZ}, omdat men via de tweede set de spati�ring kan
be�nvloeden: \type{{\,O}} in plaats van \type{O}, maar mooi is anders. 

We laten nog een voorbeeld zien. Ook hier leiden meer wegen naar Rome. Bij
een keuze voor een methode dient men de consistentie van weergeven door een
document heen te bewaken. 

\startbuffer
\startchemie
  [breedte=passend,hoogte=passend,kader=aan,
   schaal=klein]
  \chemie
    [ONE,SB15,DB7,Z057,3OFF1,MOV1,Z0,3OFF1,MOV1,
     Z017,SB1357,MOV3,Z0,MOV3,SB1357,Z013,3OFF5,
     MOV5,Z0,3OFF5,SB5,Z5]
    [C,H_2N,NH,(CH_2)_3,C,COOH,H,\SL{NH},C,COOH,H,
     (CH_2)_2,HOOC]
\stopchemie
\stopbuffer

\startfiguurtekst
  [links][]
  {geen}
  {\haalbuffer}
\voorbeeld
\typebuffer
\stopfiguurtekst

\startbuffer
\startchemie
  [breedte=passend,hoogte=passend,kader=aan,
   schaal=klein]
  \chemie [ONE,SB15,DB7,Z057,3OFF1][C,H_2N,NH]
  \chemie [MOV1,Z0,3OFF1]          [(CH_2)_3]
  \chemie [MOV1,Z017,SB1357]       [C,COOH,H]
  \chemie [MOV3,Z0]                [\SL{NH}]
  \chemie [MOV3,SB1357,Z013,3OFF5] [C,COOH,H]
  \chemie [MOV5,Z0,3OFF5,SB5,Z5]   [(CH_2)_2,HOOC]
\stopchemie
\stopbuffer

\startfiguurtekst
  [links][]
  {geen}
  {\haalbuffer}
\voorbeeld
\typebuffer
\stopfiguurtekst

\startbuffer
\startchemie
  [breedte=passend,hoogte=passend,kader=aan,
   schaal=klein]
  \chemie
    [ONE,Z0,SAVE,MOV7,SB1357,Z017,3OFF5,MOV5,Z0,
     3OFF5,MOV5,SB15,DB7,Z057,RESTORE,
     MOV3,SB1357,Z013,MOV5,3OFF5,Z0,6OFF5,SB5,Z5]
    [\SL{NH},C,COOH,H,(CH_2)_3,C,H_2N,NH,C,COOH,H,
     (CH_2)_2,HOOC]
\stopchemie
\stopbuffer

\startfiguurtekst
  [links][]
  {geen}
  {\haalbuffer}
\voorbeeld
\typebuffer
\stopfiguurtekst

\startbuffer
\startchemie
  [breedte=passend,hoogte=passend,kader=aan,
   schaal=klein]
  \chemie
    [ONE,Z0,MOV7,SB1357,Z017,3OFF5,MOV5,Z0,
     3OFF5,MOV5,SB15,DB7,Z057,MOV0,MOV3,SB1357,
     Z013,MOV5,3OFF5,Z0,6OFF5,SB5,Z5]
    [\SL{NH},C,COOH,H,(CH_2)_3,C,H_2H,NH,C,COOH,H,
     (CH_2)_2,HOOC]
\stopchemie
\stopbuffer

\startfiguurtekst
  [links][]
  {geen}
  {\haalbuffer}
\voorbeeld
\typebuffer
\stopfiguurtekst

Waar de voorkeur naar uitgaat hangt af van de smaak van programmeren en
natuurlijk de gewenste typografische kwaliteit. Let op het gebruik van 
\type{SAVE} en \type{RESTORE}. Met deze commando's kan de huidige toestand 
worden vestgelegd en opgeroepen. 

Tot slot laten we nog een voorbeeld zien van een combinatie van \type{SIX} en
\type{FIVE}. Let op het gebruik van \type{SS}. 

\startbuffer
\startchemie
    [breedte=passend,hoogte=passend,kader=aan]
  \chemie
    [SIX,DB135,SB246,Z,SR6,RZ6][C,C,N,\SR{HC},N,C,NH_2]
  \chemie
    [SIX,MOV1,DB1,SB23,SS6,Z1..5,SR3,RZ3][N,\SL{CH},N,C,C,H]
\stopchemie
\stopbuffer

\startfiguurtekst
  [links][]
  {geen}
  {\haalbuffer}
\voorbeeld
\typebuffer
\stopfiguurtekst

\hoofdstuk{Versieren}

We kunnen structuurformules voorzien van symbolen en teksten. Neem
bijvoorbeeld: 

\startbuffer
\startchemie
  [hoogte=4500,boven=1250,breedte=passend,kader=aan]
  \ondertekst
    {2,2-Dimethyl-3-etylpentan}
  \chemie
    [ONE,Z3570,SB1357]
    [CH_3,\T{1}{H_3C},CH_3,\SR{\LT{2}{C}}]
  \chemie
    [MOV1,OFF1,Z0,SB3]
    [\T{3}{CH}]
  \chemie
    [MOV3,Z0,SB3,MOV3,Z0,MOV7,MOV7]
    [CH_2,CH_3]
  \chemie
    [OFF1,SB1,MOV1,OFF1,Z0,2OFF1,SB1,Z1]
    [\T{4}{CH_2},\T{5}{CH_3}]
\stopchemie
\stopbuffer

\startfiguurtekst
  [links][]
  {geen}
  {\haalbuffer}
\voorbeeld
\typebuffer
\stopfiguurtekst

Er is een hele reeks commando's als \type{\T}. Hierbij zijn in een aantal
gevallen de argumenten optioneel. Zo kunnen we ladingsgetallen in romeinse
cijfers plaatsen met behulp van \type{\+} en \type{\-} of direct
met~\type{\1} tot en met~\type{\7}. 

\plaatstabel
{Versieringen: ladinggetallen}
\starttabel[ | l | l | ]
\HL
\VL \type{\+{nummer}} \VL positief romeins ladinggetal         \VL\FR
\VL \type{\-{nummer}} \VL negatief romeins ladinggetal         \VL\MR
\VL \type{\1}         \VL het romeinse cijfer 1 (zonder teken) \VL\MR
\VL \type{\7}         \VL idem voor 2, 3, 4, 5, 6 en 7         \VL\LR
\HL
\stoptabel

Een ladinggetal wordt gecentreerd boven het atoom waar het betrekking op
heeft. Zo resulteert: 

\startbuffer
\plaatsformule
  \startformule
    \chemie{S}+\chemie{O_2}
    \chemie{GIVES}
    \chemie{\+{4}{S}\-{2}{O_2}}
  \stopformule
\stopbuffer

\typebuffer

in:

\haalbuffer

We kunnen (eindeloze) reeksen weergeven met \type{\[} en~\type{\]}. Hierbij
zijn beide argumenten optioneel, zoals bijvoorbeeld blijkt in de onderstaande
weergave van \kap{PTFE} of Polytetrafluorethaan, beter bekend onder de naam
Teflon. 

\startbuffer
\startchemie[breedte=passend,kader=aan]
  \chemie
    [ONE,ZT5,SB5,OFF1,Z0,OFF1,SB1,MOV1,SB5,OFF1,Z0,OFF1,SB1,ZT1]
    [\[,CF_2,CF_2,\]{\sl n}]
\stopchemie
\stopbuffer

\startfiguurtekst
  [links][]
  {geen}
  {\haalbuffer}
\voorbeeld
\typebuffer
\stopfiguurtekst

\plaatstabel
{Versieringen: herhalen}
\starttabel[ | l | l | l | ]
\HL
\VL \type{\[{onder}} \VL \type{\[{boven}{onder}} \VL rechter herhalingsteken \VL\FR
\VL \type{\]{onder}} \VL \type{\]{boven}{onder}} \VL linker  herhalingsteken \VL\LR
\HL
\stoptabel

Natuurlijk kunnen we teksten links, rechts, boven of onder plaatsen. We
kunnen de commando's \type{\L}, \type{\R}, \type{\T} en \type{\B} laten
voorafgaan door \type{\X} als we de afstand wat kleiner willen hebben dan
normaal. 

\plaatstabel
{Versieringen: aan de kanten}
\starttabel[ | l | l | ]
\HL
\VL \type{\L{tekst}} \VL tekst links  plaatsen \VL\FR
\VL \type{\R{tekst}} \VL tekst rechts plaatsen \VL\MR
\VL \type{\T{tekst}} \VL tekst boven  plaatsen \VL\MR
\VL \type{\B{tekst}} \VL tekst onder  plaatsen \VL\LR
\HL
\stoptabel

Naast deze vier commando's kunnen we ook alle voor de hand liggende
combinaties gebruiken, waaronder een centreer optie.  

\leavevmode\hbox to \hsize
  {\def\sample#1{\chemie{#1{\oplus}{\ruledhbox{\ttx\string#1}}}}%
   \sample\TL\hss\sample\L \hss\sample\LC\hss\sample\BL\hss
   \sample\TR\hss\sample\R \hss\sample\RC\hss\sample\BR\hss
   \sample\LT\hss\sample\T \hss\sample\RT\hss\sample\LB\hss
   \sample\B \hss\sample\RB}

\leavevmode\hbox to \hsize
  {\def\sample#1{\chemie{\X#1{\oplus}{\ruledhbox{\ttx\string\X\string#1}}}}%
   \sample\TL\hss\sample\L \hss\sample\LC\hss\sample\BL\hss
   \sample\TR\hss\sample\R \hss\sample\RC\hss\sample\BR\hss
   \sample\LT\hss\sample\T \hss\sample\RT\hss\sample\LB\hss
   \sample\B \hss\sample\RB}

Weer een ander geval is wat we in \TEX\ terminologie {\em smashed} tekst
noemen. 

\plaatstabel
{Versieringen: smashed tekst}
\starttabel[ | l | l | ]
\HL
\VL \type{\SL{tekst}} \VL links uitlijnen  \VL\FR
\VL \type{\SM{tekst}} \VL midden uitlijnen \VL\MR
\VL \type{\SR{tekst}} \VL rechts uitlijnen \VL\LR
\HL
\stoptabel

De midden variant is hieronder weergegeven, de andere laten zich hieruit
afleiden. De tekst wordt wel gecentreerd rond het eerste karakter. 

\startbuffer
\startchemie[kader=aan]
  \chemie[SIX,B,R36,RZ36][\SL{COOH},\SL{COOH}]
\stopchemie
\stopbuffer

\startfiguurtekst
  [links][]
  {geen}
  {\haalbuffer}
\voorbeeld
\typebuffer
\stopfiguurtekst

We kunnen teksten boven, onder of midden in structuren plaatsen met de
commando's \type{\boventekst}, \type{\middentekst} en \type{\ondertekst}. De
exacte positie wordt bepaald door de hoogte en diepte van de
structuurformule. 

\startbuffer
\plaatsformule
  \startformule
    \startchemie[breedte=passend]
      \chemie[SIX,B,Z1,MOV1,B][\hbox{$\bullet$}]
      \boventekst{{\sl trans}-Decalin}
    \stopchemie
    \hskip 24pt
    \startchemie[breedte=passend]
      \chemie[SIX,B,Z12,MOV1,B][\hbox{$\bullet$},\hbox{$\bullet$}]
      \ondertekst{{\sl cis}-Decalin}
    \stopchemie
  \stopformule
\stopbuffer

\typebuffer

Beide {\em Decalin} formules zien er als volgt uit: 

\haalbuffer 

\hoofdstuk{Assenstelsel}

Structuren worden gezet in een afgeperkte ruimte, gemakshalve aangeduid als
assenstelsel. De afmetingen van dit stelsel en de plaats van de oorsprong
zijn in te stellen. Bovendien kan het assenstelsel, ten behoeve van het
positioneren in de tekst, zichtbaar worden gemaakt en kan een kader worden
getrokken. 

\startbuffer
\startchemie
  [assenstelsel=aan,
   breedte=6000,hoogte=4000]
\stopchemie
\stopbuffer

\startfiguurtekst
  [links][]{geen}
  {\haalbuffer}
\voorbeeld[vb:assenstelsel]
\typebuffer
\stopfiguurtekst

\startbuffer
\startchemie
  [assenstelsel=aan,
   links=2000,rechts=4000]
\stopchemie
\stopbuffer

\startfiguurtekst
  [links][]{geen}
  {\haalbuffer}
\voorbeeld
\typebuffer
\stopfiguurtekst

\startbuffer
\startchemie
  [assenstelsel=aan,
   breedte=6000,rechts=1000]
\stopchemie
\stopbuffer

\startfiguurtekst
  [links][]{geen}
  {\haalbuffer}
\voorbeeld
\typebuffer
\stopfiguurtekst

\startbuffer
\startchemie
  [assenstelsel=aan,
   breedte=6000,boven=1000,onder=3000]
\stopchemie
\stopbuffer

\startfiguurtekst
  [links][]{geen}
  {\haalbuffer}
\voorbeeld
\typebuffer
\stopfiguurtekst

De afmetingen van het assenstelsel bepalen de afmeting van de totale
structuur. Als echter \type{breedte=passend} en/of \type{hoogte=passend}
wordt meegegeven, dan worden de afmetingen van de totale structuur bepaald
door de werkelijke afmetingen. Waar voor gekozen wordt, hangt mede af van de
manier waarop structuren in de tekst worden geplaatst: los van elkaar, naast
elkaar, onder elkaar enz. \in{Voorbeeld}[vb:assenstelsel] toont de
standaardinstellingen. 

Binnen een \type{\start}||\type{\stop}||paar kunnen \PICTEX||macro's worden
gebruikt. Enige voorzichtigheid is daarbij natuurlijk wel geboden. 

\hoofdstuk{Instellingen}

Achter \type{\startchemie} en \type{\stelchemiein} kunnen verschillende
instellingen worden meegegeven. 

\plaatstabel
{Instellingen bij structuurformules.}
\starttabel[|lT|Tl|lT|]
\HL
\VL \bf variabele \VL \bf instellingen           \VL \bf default \VL\SR
\HL
\VL breedte       \VL getal                      \VL 4000        \VL\FR
\VL hoogte        \VL getal                      \VL 4000        \VL\MR
\VL links         \VL getal                      \VL             \VL\MR
\VL rechts        \VL getal                      \VL             \VL\MR
\VL boven         \VL getal                      \VL             \VL\MR
\VL onder         \VL getal                      \VL             \VL\LR
\HL
\VL resolutie     \VL getal                      \VL \type{\outputresolution} \VL\FR
\VL korps         \VL 8pt 9pt 10pt enz.          \VL \type{\bodyfontsize}        \VL\MR
\VL letter        \VL \type{\rm} \type{\bf} enz. \VL \type{\rm}               \VL\LR
\HL
\VL schaal        \VL getal klein middel groot   \VL middel      \VL\FR
\VL formaat       \VL klein middel groot         \VL middel      \VL\LR
\HL
\VL status        \VL start stop                 \VL start       \VL\FR
\VL optie         \VL test                       \VL             \VL\LR
\HL
\VL assenstelsel  \VL aan uit                    \VL uit         \VL\FR
\VL kader         \VL aan uit                    \VL uit         \VL\LR
\HL
\VL variant       \VL 1 2                        \VL 1           \VL\SR
\HL
\VL offset        \VL HIGH LOW                   \VL LOW         \VL\SR
\HL
\stoptabel

Standaard loopt het assenstelsel van $-2000$ tot $+2000$, zowel in de hoogte
als in de breedte. Het punt \type{Z0} ligt daarbij op (0,0). Andere
verdelingen kunnen worden ingesteld met behulp van \type{links},
\type{rechts}, \type{boven} en/of \type{onder} in combinatie met
\type{breedte} en \type{hoogte}. 

Met \type{formaat} kan de afmeting van de karakters worden ingesteld. Er
wordt daarbij achter de schermen gebruik gemaakt van de \TEX||primitieven
\type{\textsize}, \type{\scriptsize} en \type{\scriptscriptsize}. Met
\type{schaal} stelt men de afmetingen van de structuur zelf in
(\type{1..1000}). De schaal wordt mede bepaald door \type{korps}. De
trefwoorden \type{klein}, \type{middel} en \type{groot} zijn op elkaar
afgestemd. 

De instelling \type{korps} wordt gebruikt voor berekeningen en heeft geen
gevolgen voor de lopende tekst. Binnen \CONTEXT\ en \LATEX\ is deze
instelling gekoppeld aan het mechanisme dat het korps instelt. 

Zowel in \TEX\ als in \CONTEXT\ kan in de wiskundige mode gebruik gemaakt
worden van commando's als \type{\rm}, \type{\bf} en \type{\sl}. Standaard
maakt \PPCHTEX\ gebruik van \type{\rm}. Met de instelling \type{letter} kan
een andere variant worden ingesteld. In \in{voorbeeld}[vb:letter] zijn de
substituenten bijvoorbeeld {\sl slanted} gezet. 

\startbuffer
\startchemie[kader=aan,letter=\sl]
  \chemie[SIX,B,R,RZ][A,B,C,D,E,F]
\stopchemie
\stopbuffer

\startfiguurtekst
  [links][]
  {geen}
  {\haalbuffer}
\voorbeeld[vb:letter]
\typebuffer
\stopfiguurtekst

De instelling \type{letter} geldt (vooralsnog) voor zowel chemische formules
in de tekst als die in een figuur. De sub- en superscripts veranderen mee,
zoals blijkt uit: {\stelchemiein[letter=\rm]\chemie{CH_4}},
{\stelchemiein[letter=\bf]\chemie{CH_4}} en
{\stelchemiein[letter=\sl]\chemie{CH_4}}, waarbij de instellingen
achtereenvolgens zijn: \type{\rm}, \type{\bf} en \type{\sl}. Italic
\type{\it} formules leiden tot grotere regelafstanden. Binnen \CONTEXT\ zijn
standaard ook bold||slanted (\type{\bs}) en bold||italic (\type{\bi})
beschikbaar. Deze commando's passen zich automatisch aan de actuele stijl
aan: {\stelchemiein[letter=\ss\sl]\chemie{CH_4}},
{\stelchemiein[letter=\rm\sl]\chemie{CH_4}},
{\stelchemiein[letter=\tt\sl]\chemie{CH_4}} enz. (\type{\ss}, \type{\rm},
\type{\tt}). 

Natuurlijk kunnen we ook direct lettertypes instellen, zoals uit het volgende
voorbeeld blijkt. 

\startbuffer
\startchemie[kader=aan]
  \chemie[SIX,B,R,RZ][\tf a,\bf b,\it c,\sl d,\bi e,\bs f]
\stopchemie
\stopbuffer

\startfiguurtekst
  [links][]
  {geen}
  {\haalbuffer}
\voorbeeld
\typebuffer
\stopfiguurtekst

Met \type{status} kan het tijdrovende rekenwerk worden kortgesloten. De
variabelen \type{kader} en \type{assenstelsel} spreken voor zich. Met
\type{optie=test} wordt om de tekst een kader getekend, zodat men kan zien
hoe wordt uitgelijnd.\voetnoot{Binnen \CONTEXT\ wordt gebruik gemaakt van de
weergave in de visuele debugger. Deze toont de baseline als een stippellijn.
De module \type{supp-vis} is overigens ook in andere systemen te gebruiken.}
Met \type{variant} stellen we de kwaliteit van de lijnen in. Standaard
gebruikt \PICTEX\ een 5~punts . om lijnen te tekenen. Als variant~2 wordt
gekozen, worden kleinere punten gebruikt en dus dunnere lijnen getekend. 

\startbuffer
\startchemie[kader=aan,optie=test,variant=2]
  \chemie[SIX,B,R,RZ][R_1,R_2,R_3,R_4,R_5,R_6]
\stopchemie
\stopbuffer

\startfiguurtekst
  [links][]
  {geen}
  {\haalbuffer}
\voorbeeld[vb:test]
\typebuffer
\stopfiguurtekst

De offset heeft betrekking op de positie van de sub- en superscripts. Bij
\type{HIGH} worden de subscripts hoog geplaatst (\chemie{HIGH,H_2O}) en bij
\type{LOW} vanzelfsprekend wat lager (\chemie{LOW,H_2O}). 

Een structuur kan in verschillende formaten worden weergegeven. Het instellen
gebeurt met \type{formaat} en \type{schaal}. 

\startbuffer
\startchemie[kader=aan,schaal=klein,formaat=klein]
  \chemie[SIX,B,R,RZ][1,2,3,4,5,6]
\stopchemie
\stopbuffer

\startfiguurtekst
  [links][]
  {geen}
  {\haalbuffer}
\voorbeeld
\typebuffer
\stopfiguurtekst

\startbuffer
\startchemie[kader=aan,schaal=middel,formaat=middel]
  \chemie[SIX,B,R,RZ][1,2,3,4,5,6]
\stopchemie
\stopbuffer

\startfiguurtekst
  [links][]
  {geen}
  {\haalbuffer}
\voorbeeld
\typebuffer
\stopfiguurtekst

\startbuffer
\startchemie[kader=aan,schaal=groot,formaat=groot]
  \chemie[SIX,B,R,RZ][1,2,3,4,5,6]
\stopchemie
\stopbuffer

\startfiguurtekst
  [links][]
  {geen}
  {\haalbuffer}
\voorbeeld
\typebuffer
\stopfiguurtekst

Eventueel kan bij \type{schaal} een getal tussen~1 en 1000 worden ingevuld.
De waarden die bij de trefwoorden \type{klein}, \type{middel} of \type{groot}
horen zijn op elkaar afgestemd. 

\hoofdstuk{Symbolen}

Er zijn enkele symbolen beschikbaar ten behoeve van het zetten van
reactievergelijkingen. In de onderstaande figuur wordt een vergelijking
getoond. Deze vergelijking is als volgt gedefinieerd: 

\startbuffer
\stelchemiein
  [formaat=klein,
   schaal=klein,
   breedte=passend,
   hoogte=5500,
   onder=1500]

\hbox
  {\startchemie
     \chemie[SIX,B,ER6,RZ6][O]
   \stopchemie
   \startchemie
     \chemie[SPACE,PLUS,SPACE]
   \stopchemie
   \startchemie
     \chemie[FIVE,ROT4,B125,+SB3,-SB4,Z4,SR4,RZ4][N,H]
   \stopchemie
   \startchemie
     \chemie[SPACE,GIVES,SPACE][?]
   \stopchemie
   \startchemie
     \chemie[SIX,B,EB6,R6,SUB4,FIVE,ROT4,B125,+SB3,-SB4,Z4][N]
   \stopchemie
   \startchemie
     \chemie[SPACE,PLUS,SPACE,CHEM][H_2O]
   \stopchemie}
\stopbuffer

\typebuffer

De \type{\hbox} is nodig om de structuren achter elkaar te zetten. De
symbolen \type{GIVES} en \type{PLUS} spreken voor zich. Met \type{SPACE} kan
extra ruimte worden afgedwongen. 

\plaatsfiguur
  [hier]
  [fig:reactie]
  {geen}
  {\haalbuffer}

Een evenwicht kan worden weergegeven met \type{EQUILIBRIUM}. Boven
\type{GIVES} en \type{EQUILIBRIUM} kan een tekst worden gezet. In het
voorbeeld is dat een '?'. 

Naast \type{GIVES} en \type{EQUILIBRIUM} is ook \type{MESOMERIC} beschikbaar.
Daarnaast kunnen complexen worden gezet met behulp van \type{OPENCOMPLEX} en
\type{CLOSECOMPLEX}. 

\hoofdstuk{Positioneren}

Bij het combineren van moleculen, bijvoorbeeld met \type{SUB}, wisselen een
aantal kenmerkende maten en posities, zoals het middelpunt. Het is daarom
mogelijk de oude toestand te bewaren en weer op te roepen met \type{SAVE} en
\type{RESTORE}. 

\plaatstabel
{Positioneren.}
\starttabel[ | r T | l w(4cm) | p(6cm) | ]
\HL
\VL SAVE    \VL Save Status    \VL het opslaan van de huidige toestand   \VL\FR
\VL RESTORE \VL Restore Status \VL het herstellen van de huidige toestand\VL\LR
\HL
\stoptabel

We gebruiken de commando's \type{SAVE} en \type{RESTORE} vooral bij
substituenten, Bij radicalen vallen we veelal terug op \type{PB} en
\type{PE}, wat in feite positioneringscommando's zijn. 

\startbuffer
\definieerchemie[molecuul]
  {\chemie
     [ONE,Z0,SB1357,
      SAVE,SUB2,SIX,B,R6,C,RESTORE,
      MOV1,Z0,SB137,
      MOV1,Z0,SB37,
      MOV1]
     [C,C,C]}

\startchemie[breedte=passend,hoogte=passend]
  \chemie[molecuul,molecuul,molecuul]
\stopchemie
\stopbuffer

\typebuffer

Dit voorbeeld toont tevens de mogelijkheid verschillende moleculen aaneen te
schakelen tot ketens. 

\plaatsfiguur
  [hier]
  [fig:save]
  {geen}
  {\haalbuffer}

Het onderstaande voorbeeld is wat geavanceerder en toont tevens de
mogelijkheid formules te zetten. Bij dergelijke formules is het instellen van
de hoogte belangrijk. 

\startbuffer
\plaatsformule
  \startformule
    \stelchemiein
      [breedte=passend,boven=2000,onder=2000,
       schaal=klein,formaat=klein]
    \startchemie
      \chemie
        [ONE,
         SAVE,
           Z0,SB7,SB3,SB1,MOV1,Z0,SB1,MOV1,Z0,DB8,CZ8,SB1,Z1,
         RESTORE,
         SAVE,
           SUB4,ONE,Z0,SB3,SB1,MOV1,Z0,SB1,MOV1,Z0,DB8,CZ8,SB1,Z1,
         RESTORE,
         SUB2,ONE,Z0,SB7,SB1,MOV1,Z0,SB1,MOV1,Z0,DB8,CZ8,SB1,Z1]
        [\SR{HC},O,C,O,C_{19}H_{39},
         \SR{H_{2}C},O,C,O,C_{17}H_{29},
         \SR{H_{2}C},O,C,O,C_{21}H_{41}]
    \stopchemie
    \startchemie
      \chemie[SPACE,PLUS,SPACE]
    \stopchemie
    \startchemie[rechts=600]
      \chemie[ONE,CZ0][3CH_{3}OH]
    \stopchemie
    \startchemie
      \chemie[SPACE,GIVES,SPACE,SPACE][H^+/H_2O]
    \stopchemie
    \startchemie
      \chemie
        [ONE,
         SAVE,
           Z0,SB7,SB3,SB1,Z1,
         RESTORE,
         SAVE,
           SUB4,ONE,Z0,SB3,SB1,Z1,
         RESTORE,
         SUB2,ONE,Z0,SB7,SB1,Z1]
        [\SR{HC},OH,
         \SR{H_{2}C},OH,
         \SR{H_{2}C},OH]
    \stopchemie
    \startchemie
      \chemie[SPACE,PLUS,SPACE]
    \stopchemie
    \startchemie
      \chemie
        [ONE,
         SAVE,
           Z0,DB8,CZ8,SB1,SB5,Z5,MOV1,Z0,SB1,Z1,
         RESTORE,
         SAVE,
           SUB4,ONE,Z0,DB8,CZ8,SB1,SB5,Z5,MOV1,Z0,SB1,Z1,
         RESTORE,
         SUB2,ONE,Z0,DB8,CZ8,SB1,SB5,Z5,MOV1,Z0,SB1,Z1]
        [C,O,C_{19}H_{39},O,CH_{3},
         C,O,C_{17}H_{29},O,CH_{3},
         C,O,C_{21}H_{41},O,CH_{3}]
    \stopchemie
  \stopformule
\stopbuffer

\typebuffer

Deze definitie zou wat compacter kunnen (\type{SB731} in plaats van
\type{SB7,SB3,SB1}) maar dit komt de leesbaarheid niet altijd ten goede.
Dergelijke omvangrijke structuren kan men overigens het best stapsgewijs
opbouwen in een aparte file (eventueel lekker snel in \PLAIN\ \TEX). 

\haalbuffer

We geven nog twee voorbeelden waarin we bovendien tekst plaatsen onder de
structuur. 

\startbuffer
\plaatsformule
  \startformule
    \stelchemiein
      [breedte=passend,boven=1500,onder=3500]
    \startchemie
      \chemie
        [ONE,Z0,DB1,SB3,SB7,Z7,MOV1,Z0,SB3,SB7,Z3,Z7,
         MOV0,SUB2,SIX,B,R6,C]
        [C,H,C,H,H]
      \ondertekst{styreen}
    \stopchemie
    \quad\quad\quad
    \startchemie
      \chemie
        [ONE,Z0,DB1,SB3,SB7,Z3,Z7,
         MOV1,Z0,SB1,SB3,Z3,
         MOV1,Z0,DB1,SB3,Z3,
         MOV1,Z0,SB3,SB7,Z3,Z7]
        [C,H,H,C,H,C,H,C,H,H]
      \ondertekst{1,3-butadieen}
   \stopchemie
  \stopformule
\stopbuffer

\typebuffer

\haalbuffer

Het gebruik van \type{OFF} kijkt soms nauw. De onderstaande voorbeelden laten
zien dat subtiele verplaatsingen mogelijk zijn. 

\startbuffer
\startchemie[hoogte=passend,kader=aan]
  \chemie[ONE,SB]
\stopchemie
\stopbuffer

\startfiguurtekst
  [links][]
  {geen}
  {\haalbuffer}
\voorbeeld
\typebuffer
\stopfiguurtekst

\startbuffer
\startchemie[hoogte=passend,kader=aan]
  \chemie[ONE,3OFF1,SB]
\stopchemie
\stopbuffer

\startfiguurtekst
  [links][]
  {geen}
  {\haalbuffer}
\voorbeeld
\typebuffer
\stopfiguurtekst

\startbuffer
\startchemie[hoogte=passend,kader=aan]
  \chemie[ONE,MOV1,SB]
\stopchemie
\stopbuffer

\startfiguurtekst
  [links][]
  {geen}
  {\haalbuffer}
\voorbeeld
\typebuffer
\stopfiguurtekst

\startbuffer
\startchemie[hoogte=passend,kader=aan]
  \chemie[ONE,3OFF1,MOV1,SB]
\stopchemie
\stopbuffer

\startfiguurtekst
  [links][]
  {geen}
  {\haalbuffer}
\voorbeeld
\typebuffer
\stopfiguurtekst

\startbuffer
\startchemie[hoogte=passend,kader=aan]
  \chemie[ONE,MOV1,3OFF1,SB]
\stopchemie
\stopbuffer

\startfiguurtekst
  [links][]
  {geen}
  {\haalbuffer}
\voorbeeld
\typebuffer
\stopfiguurtekst

\startbuffer
\startchemie[hoogte=passend,kader=aan]
  \chemie[ONE,MOV1,3OFF1,OFF0,SB]
\stopchemie
\stopbuffer

\startfiguurtekst
  [links][]
  {geen}
  {\haalbuffer}
\voorbeeld
\typebuffer
\stopfiguurtekst

\startbuffer
\startchemie[hoogte=passend,kader=aan]
  \chemie[ONE,MOV1,3OFF1,MOV0,SB] 
\stopchemie
\stopbuffer

\startfiguurtekst
  [links][]
  {geen}
  {\haalbuffer}
\voorbeeld
\typebuffer
\stopfiguurtekst

\startbuffer
\startchemie[hoogte=passend,kader=aan]
  \chemie[ONE,MOV1,MOV0,SB]
\stopchemie
\stopbuffer

\startfiguurtekst
  [links][]
  {geen}
  {\haalbuffer}
\voorbeeld
\typebuffer
\stopfiguurtekst

Het volgende voorbeeld laat zien hoe complexen kunnen worden weergegeven. Let
op het gebruik van \type{RBT}. Normaal is de extra spati�ring niet nodig,
maar hier gebruiken we |<|het commando is niet zichtbaar|>| een nog kleiner
lettertype omdat de figuur anders te breed wordt. 

\startbuffer
\startformule
\stelchemiein[schaal=klein,breedte=passend]
\startchemie
  \chemie[OPENCOMPLEX,SPACE]
\stopchemie
\startchemie
  \chemie[SIX,B,EB35,R6,-R1,+R1]                  
  \chemie[SIX,PB:RZ6,ONE,OFF1,Z0,EP57,PE][\SL{OH}]
  \chemie[SIX,-RZ1,+RZ1,RT2][H,Br,\oplus]
\stopchemie
\startchemie
  \chemie[SPACE,MESOMERIC,SPACE]
\stopchemie
\startchemie
  \chemie[SIX,B,EB25,R6,-R1,+R1]                  
  \chemie[SIX,PB:RZ6,ONE,OFF1,Z0,EP57,PE][\SL{OH}] 
  \chemie[SIX,-RZ1,+RZ1,RT4][H,Br,\oplus]
\stopchemie
\startchemie
  \chemie[SPACE,MESOMERIC,SPACE]
\stopchemie
\startchemie
  \chemie[SIX,B,EB24,R6,-R1,+R1]
  \chemie[SIX,PB:RZ6,ONE,OFF1,Z0,EP57,PE][\SL{OH}]  
  \chemie[SIX,-RZ1,+RZ1,RBT6][H,Br,\ \oplus]
\stopchemie
\startchemie
  \chemie[SPACE,MESOMERIC,SPACE]
\stopchemie
\startchemie
  \chemie[SIX,B,EB24,ER6,-R1,+R1]                   
  \chemie[SIX,PB:RZ6,ONE,OFF1,Z0,EP5,ZT7,PE][\SL{OH},\oplus]
  \chemie[SIX,-RZ1,+RZ1][H,Br]
\stopchemie
\startchemie
  \chemie[SPACE,CLOSECOMPLEX]
\stopchemie
\stopformule
\stopbuffer

\typebuffer

Zonder het gebruik van \type{SPACE} zouden de verschillende afbeeldingen te
dicht op elkaar komen te staan. Het optimaliseren van het afbeelden van
dergelijke reactievergelijkingen is meestal een iteratief proces. 

\start
\switchnaarkorps[klein]
\haalbuffer
\stop

\hoofdstuk{Lopende tekst}

Naast het zetten van structuurformules wordt ook het zetten van
reactievergelijkingen ondersteund. Het hiervoor beschreven commando
\type{\chemie} heeft daarom nog twee uitvoeringen: 

\starttypen
\chemie{formule}
\chemie{formule}{onder tekst}
\chemie{formule}{boven tekst}{onder tekst}
\stoptypen

Dit commando past zich aan de plaats in de tekst aan. Dat wil zeggen dat
onderscheid wordt gemaakt tussen: 

\startopsomming[opelkaar]
\som tekst||mode
\som wiskundige tekst||mode
\som wiskundige display||mode
\stopopsomming

Als het commando in de lopende tekst wordt gegeven, dan worden automatisch
\type{$ $} om het commando geplaatst. Zo levert \type{\chemie{NH_4^+}} de
formule \chemie{NH_4^+} op. 

Hetzelfde resultaat wordt bereikt als het commando tussen \type{$ $} wordt
geplaatst. In beide gevallen wordt het tweede en derde argument dus
weggelaten. Als we het commando tussen \type{$$ $$} (of \type{\startformule}
en \type{\stopformule}) plaatsen, dan hebben de extra (maar facultatieve)
argumenten wel een functie. Eerst laten we de eenvoudige variant zien. 

\startbuffer
\plaatsformule
  \startformule
    \chemie{2H_2} \chemie{PLUS} \chemie{O_2}
    \chemie{GIVES} \chemie{2H_2O}
  \stopformule
\stopbuffer

\typebuffer

levert:

\haalbuffer

Het chemische deel kan ook korter worden gespecificeerd: 

\starttypen
\chemie{2H_2,PLUS,O_2,GIVES,2H_2O}
\stoptypen

of zelfs:

\starttypen
\chemie{2H_2,+,O_2,->,2H_2O}
\stoptypen

Het zal de \TEX||kenner opvallen dat zowel het plus||teken als de pijl op de
baseline staan. Vergelijk $+$ maar eens met \chemie{+}. Er is namelijk voor
gezorgd dat, onafhankelijk van het formaat van weergeven, de \chemie{+} en de
\chemie{->} in elkaars verlengde liggen. 

Naast \type{PLUS} en \type{GIVES} kan \type{EQUILIBRIUM} (\type{<->}) worden
opgegeven, wat de dubbele pijlen \chemie{EQUILIBRIUM} oplevert. Met
\type{MESOMERIC} of \type{<>} krijgt men \chemie{MESOMERIC}. 

Deze formule kan ook in de lopende tekst worden opgenomen. In dat geval wordt
een wat minder ruime layout gekozen: \chemie{2H_2,+,O_2,->,2H_2O}. Het is
daarnaast mogelijk bindingen weer te geven. Zo levert
\type{\chemie{H,SINGLE,CH,DOUBLE,HC,SINGLE,H}} de chemische formule
\chemie{H,SINGLE,CH,DOUBLE,HC,SINGLE,H}. Ook dit kan korter:
\type{\chemie{H,-,CH,--,HC,-,H}}. Een drievoudige binding is op te roepen met
\type{TRIPLE} of \type{---}: \chemie{HC,TRIPLE,CH}. 

We keren nog even terug naar de display||mode. Het tweede en derde argument
kunnen worden gebruikt om een toelichting op de formule te geven: 

\startbuffer
\plaatsformule
  \startformule
    \chemie{2H_2}{waterstof} \chemie{PLUS} \chemie{O_2}{zuurstof}
    \chemie{GIVES}{heftig} \chemie{2H_2O}{water}
   \stopformule
\stopbuffer

\typebuffer

We kunnen dus ook onder (en boven) symbolen plaatsen! 

\haalbuffer

Er zijn maximaal drie argumenten mogelijk, waarbij het laatste argument onder
wordt geplaatst. 

\startbuffer
\plaatsformule
  \startformule
    \chemie{H_2O}{liquid}{water}
    \hbox{c.q.}
    \chemie{H_2O}{water}
  \stopformule
\stopbuffer

\haalbuffer

De bovenstaande formule is gezet met:

\typebuffer

Het formaat van de in de tekst opgenomen formules kan worden ingesteld met
het setup||commando: 

\plaatstabel
{Instellingen bij tekstformules.}
\starttabel[|lT|Tl|lT|]
\HL
\VL \bf variabele \VL \bf instellingen   \VL \bf default \VL\SR
\HL
\VL tekstformaat  \VL klein middel groot \VL groot       \VL\SR
\HL
\stoptabel

Het commando \type{\chemie{H,SINGLE,CH,DOUBLE,HC,SINGLE,H}} levert bij
achtereenvolgens de instellingen \type{groot}, \type{middel} en \type{klein}
de formules: 

\leavevmode\hbox
  {\stelchemiein[tekstformaat=groot]\chemie{H,-,CH,--,HC,-,H}\hskip2em
   \stelchemiein[tekstformaat=middel]\chemie{H,-,CH,--,HC,-,H}\hskip2em
   \stelchemiein[tekstformaat=klein]\chemie{H,-,CH,--,HC,-,H}}

\hoofdstuk{Subscripts}

Sub- en superscripts worden, zoals door Knuth in het {\TeX}Book wordt
aanbevolen, wat lager geplaatst. Zo is het in het {\TeX}Book op pagina~179
gegeven, chemisch gezien wat vreemde, voorbeeld te zetten door in de tekst
\type{\chemie{Fe_2^{+2}Cr_2O_4}} op te nemen. Dit levert
\chemie{Fe_2^{+2}Cr_2O_4}. Zonder correctie zou dit zijn geweest: $\rm
Fe_2^{+2}Cr_2O_4$. 

De plaats van het subscript wordt standaard bepaald door de instelling van
\type{offset}: \type{HIGH} of \type{LOW}. Deze instelling kan worden
overruled door de gelijknamige commando's. 

\startbuffer
\startchemie[kader=aan,optie=test,variant=2]
  \chemie[SIX,B,R13,HIGH,RZ1,LOW,RZ3][NH_3,NH_3]
\stopchemie
\stopbuffer

\startfiguurtekst
  [links][]
  {geen}
  {\haalbuffer}
\voorbeeld[vb:test]
\typebuffer
\stopfiguurtekst

Hoewel dus per substituent de plaats van het subscript kan worden ingesteld,
verdient het aanbeveling deze instelling globaal te doen, zodat de formules
onderling consistent worden gezet. 

De keywords \type{LOW} en \type{HIGH} kunnen ook worden gebruikt in
tekstformules, zo resulteert \type{\chemie{HIGH,H_2O}} in \chemie{HIGH,H_2O}
en \type{\chemie{LOW,H_2O}} in \chemie{LOW,H_2O}. 

\vervolgblad
  {Achtergronden}
  [links=1500,breedte=4000]
  [FIVE,ROT3,B125,+SB3,-SB4,ER35,SR4,CRZ35,Z4,RZ4,Z0]
  [O,O,N,CH_3]

\hoofdstuk{Installatie}

Het pakket \PPCHTEX\ is in eerste instantie ontwikkeld voor gebruik binnen
\CONTEXT. Binnen dit macropakket kan \PPCHTEX\ worden geactiveerd met het
commando:\voetnoot{De eigenlijke macro's staan in de file \type{ppchtex.tex},
die wordt geladen in \type{m-chemie.tex} (de~\type{m} staat hier voor
module).} 

\starttypen
\gebruikmodules[pictex,chemie]
\stoptypen

Dit commando laad eerst de \PICTEX\ macros zodat \PPCHTEX\ weet welke output
gewenst is. De chemie macro's zijn dus niet standaard beschikbaar. 

Daarnaast kan het pakket worden gebruikt in combinatie met \PLAIN\ \TEX\ of
\LATEX. In dat geval wordt tevens de file \type{m-chemie.sty} gebruikt. Het
activeren van \PPCHTEX\ vindt in dat geval plaats door middel van
\type{\documentstyle}: 

\starttypen
\documentstyle[m-pictex,m-chemie]{}
\stoptypen

In \LATEX$2_{\textstyle\varepsilon}$ ziet het er wat anders uit: 

\starttypen
\usepackage{m-pictex}
\usepackage{m-chemie}
\stoptypen

De file \type{m-pictex} zorgt er voor dat \PICTEX\ zo effici�nt mogelijk
wordt geladen. Dit is nodig omdat \LATEX, zeker als er wat style files worden
geladen, nogal wat \type{\dimen}'s alloceert. Overigens is als gevolg van de
gebruikersinterface een wat grotere versie van \TEX\ nodig dan de standaard
implementatie. Meestal is deze wel voorradig. 

\PPCHTEX\ kan in verschillende talen worden toegesproken. De actuele taal
hangt af van de wijze van laden: 

\startregelcorrectie
\starttabel[|l|l|]
\HL 
\NC \bf versie             \NC \bf files                             \NC\SR 
\HL 
\NC nederlandtalige versie \NC \type{m-che-nl} = \type{m-chemie.sty} \NC\FR
\NC engelstalige versie    \NC \type{m-che-en} = \type{m-chemic.sty} \NC\MR
\NC duitstalige versie     \NC \type{m-che-de}                       \NC\LR
\HL 
\stoptabel
\stopregelcorrectie

In de file \type{ppchtex.noc} is te zien hoe buiten \CONTEXT\ de koppeling
met het omhullende macropakket (kan) worden gelegd. Deze file laad bovendien
een van de systeem modules van \CONTEXT\ en sluit wat zaken kort die niet
beschikbaar zijn buiten dit pakket. 

De totale distributie bestaat dus uit de definitie||files \type{ppchtex.tex}
en \type{ppchtex.noc} de opstart||files \type{m-che-nl.tex},
\type{m-che-en.tex} en \type{m-che-de.tex} en de alternatieve opstart||files
\type{m-chemie.tex} en \type{m-chemic.tex}. 

\hoofdstuk{Uitbreidbaarheid}

Het staat de gebruikers van \PPCHTEX\ natuurlijk vrij de macro's die eraan
ten grondslag liggen ook op een andere wijze in te zetten. Enige
voorzichtigheid is echter geboden omdat de macro's nog steeds worden
uitgebreid, worden geoptimaliseerd en meer robuust worden gemaakt. Sommige
macro's lijken misschien nodeloos ingewikkeld, maar schijn bedriegt.
Commando's met de vorm \type{\stel...in} maken bijvoorbeeld gebruik van
macro's die nesting en verschillende \ASCII||layouts ondersteunen. Vergelijk
bijvoorbeeld: 

\start
\steltypenin[optie=kleur]
\starttypen
\stelchemiein[formaat=klein]
\stoptypen
\stop

met:

\start
\steltypenin[optie=kleur]
\starttypen
\stelchemiein
  [formaat=klein,
   schaal=500,
   tekstformaat=groot]
\stoptypen
\stop

De instellingen mogen in een willekeurige volgorde worden opgegeven. Waar
mogelijk worden spaties en regelovergangen onderdrukt en worden meldingen
gegenereerd met betrekking tot fouten. 

Omdat \PPCHTEX\ in eerste instantie is bedoeld (en ontwikkeld) als een module
bij \CONTEXT, verstaat het pakket meerdere talen. Op dit moment worden een
nederlands-, engels- en duitstalige interface ondersteund. 

Een blik in de file \type{ppchtex.tex} leert dat bij het interpreteren van de
in \type{\chemie} tussen \type{[ ]} opgegeven commando's, gebruik wordt
gemaakt van \type{\processaction}||macro's. Deze macro's zijn, omdat ook hier
zowel nesting als een overzichtelijke layout mogelijk moet zijn, relatief
traag. Daar staat tegenover dat de \PPCHTEX||macro's inzichtelijk blijven.
Inmiddels zijn zowel de macro's zelf als het gebruik van de macro's binnen
\PPCHTEX\ redelijk geoptimaliseerd. Testen met de stopwatch hebben geleerd
dat verdere optimalisatie nauwelijks tijdwinst oplevert en wel ten koste gaat
van de flexibiliteit. 

Bij het ontwikkelen van \PPCHTEX\ is rekening gehouden met de wijze waarop
commando's binnen de produktie||omgeving \TEXEDIT\ worden weergegeven. In dit
programma worden \TEX||commando's en symbolen met een speciale functie in
kleur weergegeven, zoals hierboven. 

\hoofdstuk{Fonts}

De macro's maken gebruik van de wiskundige mode en dientengevolge van
\type{\textfont}, \type{\scriptfont} en \type{\scriptscriptfont}. Wanneer dat
nodig is, worden de \type{\fontdimen}'s~14, 16 en~17 van \type{\font2}
aangepast. Daarbij wordt rekening gehouden met de actuele korpsgrootte, die
wordt afgeleid uit de $x$-height (\type{\fontdimen5}). 

Wijzigingen in de \type{\fontdimen}'s hebben een globaal karakter, zodat
groeperen geen zin heeft. Vandaar dat de dimensies steeds worden geset en
gereset. De huidige oplossing is misschien niet de beste, maar andere
oplossingen geven problemen met geschaalde fonts. 

Het plaatsen van de atomen en moleculen (tekst) kost relatief veel
procestijd. Deze snelheid hangt mede af van de complexiteit van de
aangeroepen macro \type{\rm}. Een meer effici�nte implementatie is mogelijk. 

\hoofdstuk{Definities}

\PPCHTEX\ sluit voor wat betreft de interactie aan bij \CONTEXT. Dit betekent
dat de interface meertalig is. Tegenover het voordeel dat een ieder in zijn
eigen taal kan communiceren met het pakket, staat het nadeel dat
gemeenschappelijk gebruik van macro's wat lastiger is. 

Op dit moment zijn veel onderliggende commando's van \CONTEXT\
nederlandstalig. Hetzelfde geldt voor de parameters. \PPCHTEX\ heeft echter
al wel engelstalige commando's. Om internationaal gebruik mogelijk te maken,
moeten bibliotheken worden gedefinieerd met engelstalige commando's,
bijvoorbeeld: 

\starttypen
\startchemical
  \chemical[SIX,B,C]
\stopchemical
\stoptypen

Instellingen zijn wat lastiger. We gebruiken daarvoor systeemconstanten en 
variabelen, die helaas op dit moment nog nederlandstalig zijn. Op niet al te 
lange termijn zullen echter ook deze engelstalig zijn. Gelukkig zijn de 
constanten en variabelen goed te herkennen, zodat aanpassen niet moeilijk 
is. 

\starttypen
\setupchemical[\c!breedte=10cm,\c!hoogte=\v!passend]
\stoptypen

Parameters worden voorafgegaan door \type{\c!} en instellingen met 
\type{\v!}. Dit gaat alleen goed als \type{!} een letter is. Daarom moeten 
dergelijke instellingen worden omringt door \type{\unprotect} en 
\type{\protect}. We krijgen dus bijvoorbeeld: 

\starttypen
\unprotect

\setupchemical[\c!breedte=10cm,\c!hoogte=\v!passend]

\startchemical
  \chemical[SIX,B,C]
\stopchemical

\protect
\stoptypen

Meer informatie over de interface en de werking ervan is te vinden in de 
documentatie van de \CONTEXT\ modules uit de \type{mult} groep. 

\hoofdstuk{Kleur}

Het is (binnen \CONTEXT) mogelijk delen van een structuur te kleuren.
In~\in{voorbeeld}[vb:kleur] wordt zowel de substituent als de binding naar de
substituent rood gekleurd. 

\startbuffer
\startchemie[assenstelsel=aan,kader=aan,breedte=5000]
  \chemie[SIX,B]
  \kleur[rood]{\chemie[SIX,R1,RZ1][COOH]}
\stopchemie
\stopbuffer

\startfiguurtekst
  [links][]
  {geen}
  {\haalbuffer}
\voorbeeld[vb:kleur]
\typebuffer
\stopfiguurtekst

Het kleur||mechanisme dient natuurlijk wel eerst te zijn geactiveerd met het
commando \type{\stelkleurenin[status=start]}. 

\hoofdstuk[txt:cooh]{Interactie}

\PPCHTEX\ ondersteunt, althans bij gebruik met \CONTEXT, het opmaken van
interactieve teksten. Onder een interactieve tekst verstaan we een tekst die
kan worden geraadpleegd op de computer, waarbij delen van die tekst actief
kunnen zijn. Dat wil zeggen dat het aanklikken van bijvoorbeeld een woord
resulteert in het gaan naar de toelichting die bij dit woord hoort. In de
tekst op \naar{\chemie{COOH}}[sub:cooh] klikken kan bijvoorbeeld als gevolg
hebben dat de cursor op de betreffende substituent in een structuur gaat
staan. Andersom kan klikken op de substituent tot gevolg hebben dat naar de
uitleg over deze substituent wordt gesprongen. 

\startbuffer
\startchemie[assenstelsel=aan,kader=aan,breedte=5000]
  \chemie[SIX,B]
  \chemie[sub:cooh][SIX,R1,RZ1][COOH]
\stopchemie
\stopbuffer

\startfiguurtekst
  [links][]
  {geen}
  {\haalbuffer}
\voorbeeld[vb:interactie]
\typebuffer
\stopfiguurtekst

We zien dat er een derde argument wordt meegegeven: de verwijzing
\type{[sub:cooh]}. Dit betekent dat we vanuit de tekst kunnen verwijzen naar
de substituent \naar{\chemie{COOH}}[sub:cooh]: 

\starttypen
... tekst ... \naar{\chemie{COOH}}[sub:cooh] ... tekst ...
\stoptypen

Hierbij is \type{\naar} een \CONTEXT||commando. Klikken op de als zodanig
gemarkeerde tekst resulteert dus in het plaatsen van de cursor bij de
substituent in de afbeelding. 

Andersom kan ook. We kunnen vanuit een structuur verwijzen naar een tekst. De
verwijzing moet natuurlijk wel bestaan. 

Klikken op \chemie{COOH} in de figuur resulteert in het springen naar een
bijbehorende tekst, bijvoorbeeld gemarkeerd met: 

\starttypen
\paragraaf[txt:cooh]{Substituenten}

... tekst ... \chemie{COOH} ... tekst ... 
\stoptypen 

Een combinatie is ook mogelijk. In dat geval is het verstandig eerst te
markeren met \type{\chemie} en vanuit de formule te verwijzen met
\type{\naarchemie}. 

De koppeling met het interactiemechanisme loopt via de macro's: 

\starttypen
\localgotochemical   {verwijzing} {tekst}
\localthisischemical {verwijzing}
\stoptypen

Voor meer uitleg verwijzen we naar de source. Deze koppeling kan ook voor
andere doeleinden worden gebruikt. Enige ervaring met \TEX\ strekt daarbij
wel tot aanbeveling. 

\vervolgblad
  {Overzichten}
  [links=1250,breedte=4000,onder=1500,hoogte=4000]
  [SIX,B,C,MOV1,B,ER6,CRZ6]
  [O]

\input man-comm.tex

\stoptekst
